% Generated by scripts/materialize_unified_paper.py; do not edit.
\section{Pinned source interface and quantifier ledger}\label{low:app:bridge}

This appendix records the exact external boundary of the lower-bound
argument.  It separates version-pinned conditional probability statements
from the explicit weighted-extension rule and imports neither the
first-moment conversion nor the Ramsey conclusion proved in
\cref{low:sec:rate-lower}.

\subsection{Content-addressed source snapshots}

The source files are Hunter--Milojevi\'c--Sudakov,
arXiv:2512.17718v2, file \texttt{arXiv\_version.tex}, SHA-256
\begin{center}
\small\ttfamily
b72958ac35554eccb94dedab5800349d2c021af7d60f767125cb46998e0fd54a
\end{center}
and Lin--Niu, arXiv:2605.25843v2, file
\texttt{off-diagonal-Ramsey-R.tex}, SHA-256
\begin{center}
\small\ttfamily
dbf1bfa2c7603c81b2e56a97028b138ccf9a9299a3d1f20c0b08cf489136f429
\end{center}
The repository manifest records the versioned retrieval URLs and independently
checks these hashes.

\subsection{Source-line map for the source interface}

\begin{longtable}{@{}p{0.25\textwidth}p{0.67\textwidth}@{}}
\toprule
Pinned source lines & Exact use in this paper\\
\midrule
\endhead
HMS 145--170, 443--498&
The Gaussian graph, Bartlett representation, column exposure, edge events,
perfect events, and the unweighted projection state.  These are the
probability-space and history definitions in (2.24)--(2.33).\\
HMS 955--971, 1192--1218&
Optimized red-perfectness and the exact cutoff approximation.  Cauchy--Schwarz gives
\(\sqrt d\,|h_{ij}|\le\alpha_R^2\ell/\sqrt d\le\omega_0\), while the
diagonal approximation contributes \(d^{-1/5}\).  This is (S1).\\
HMS 330--354, 878--935&
The subgaussian definition, square-MGF lemma, pre-Cauchy centered-linear
allocation, and centered quadratic estimate.  These give the quantified
statements (2.35)--(2.36), the uniform \(t_0\) factor in (5.12), and the
square estimate in (5.14); no final Cauchy coarsening is imported.\\
Lin--Niu 359--493, 499--537&
Monotonicity of upper-truncated Gaussian variance and the exact
positive-direction CGF.  After replacing their upper cutoff by \(-b\),
these give (2.4)--(2.6) and \cref{low:lem:mills}.\\
HMS 1192--1289&
The red law of total probability, cutoff conditioning, projection merge, and
scalar reverse-induction allocation.  Equations (2.33), (2.39), and (2.40)
record the algebra that the separately assumed weighted extension (S3) must
preserve; (2.41) is the frozen scalar baseline.  The nonuniform estimate is
proved locally and is not attributed verbatim to HMS.\\
HMS 385--410 and 955--971&
For every specified retained set, the red-perfectness failure factor is
\((p/10)^{10\ell u}\).  This is the sole external extraction input (2.42);
the restoration estimates and binomial sum are proved in (6.5)--(6.11).\\
HMS 977--1181, 1294--1312, together with 385--410&
The unchanged blue reverse induction, its perfect bound, and the same greedy
extraction argument, yielding the raw probability inequality (2.43).  The
first-moment threshold (7.4) is derived locally.\\
\bottomrule
\end{longtable}

The separately fetched Yang--Mao v1 member has SHA-256
\begin{center}
\small\ttfamily
155b7104ec5b6935a576ae9f2b161976a966b0b46bd2b69153c0934ca688da2a
\end{center}
agreeing with the repository manifest.  It is not an input to the lower-bound
component, so no Yang--Mao line range is invoked here; assigning one would
misstate the lower component's dependency boundary.

\subsection{Why the cutoff window is
\texorpdfstring{\(\omega_0\)}{omega-0}}

On the event \(B^R_r\),
\[
 |b_{ij}-c|\le
 \sqrt d\,|h_{ij}|+d^{-1/5}
 \le\frac{\alpha_R^2\ell}{\sqrt d}+d^{-1/5}
 \le\omega_0+d^{-1/5}.
\tag{L-B.B.1}\label{low:app-B:eq:b-1}
\]
For every fixed \(w>\omega_0\), choosing
\(\ell\ge\ell_0(C,w)\) makes the last expression at most \(w\).
The proof is first carried out for this open range of \(w\), and only after
the fixed-\(C\), \(\ell\to\infty\) limit is taken do we let
\(w\downarrow\omega_0\).  No finite-\(\ell\) history is claimed to lie in
the closed boundary window.  This also explains why \(p\), rather than
\(p_C\), occurs in \(\alpha_R\) and \(\omega_0\).

\subsection{The old exact-CGF entry is an input formula, not a conclusion}

At a step with \(k\) future vertices, define
\[
 b_j=b_{ij},\qquad
 M_j=\sum_{q\ne j}m(b_q),\qquad
 u_j^0=\frac{m_0M_j}{\sqrt d}.
\tag{L-B.B.2}\label{low:app-B:eq:b-2}
\]
Stopping the HMS centered-moment calculation before its final Cauchy
coarsening gives the allocation \(\sum_j(u_j^0)^2\).  Applying the exact
positive-direction CGF with \(P^0_{k,d}\) gives
\[
 \sum_j\frac1{P^0_{k,d}}K_{b_j}(P^0_{k,d}u_j^0).
\tag{L-B.B.3}\label{low:app-B:eq:b-3}
\]
Their difference is the same-history quantity \(\mathfrak d_i(b)\) in
(5.16).  The cutoff monotonicities give (5.17), and summing its deterministic
lower bound gives (2.15).  With \(k/\ell\to x\) and
\(\sqrt d/\ell\to D\), the corresponding Riemann sum is (2.17).  Hence
\(G_*(C)\) is defined and charged before the controlled residual is added;
it is not imported through a final Ramsey theorem.

\subsection{Raw ledgers and the local crossing calculation}

The red interface before the new two deficits is the explicit envelope
\[
 \mathcal H_C(r)=p^{\binom r2}
 \exp\left[-\vartheta_C\frac{a^3}{p^3\sqrt d}\binom r3
 +\frac apd^{-1/5}\binom r2+r\right].
\tag{L-B.B.4}\label{low:app-B:eq:b-4}
\]
The blue interface is the explicit probability estimate (2.43).  Dividing
their logarithms by \(\ell^2\), and only then applying
\(\binom nk\le n^k\), gives the red and blue thresholds (7.2) and (7.4).
The comparison is not assumed: the source asymptotics imply
\[
 \frac{B_R(C)}2=\frac1{12}+o(1),\qquad
 \frac{CB_B(C)}2=\frac12+o(1),
\tag{L-B.B.5}\label{low:app-B:eq:b-5}
\]
whereas
\[
 G_*(C)+\widehat H_*(C)=O(1/\log C)=o(1).
\tag{L-B.B.6}\label{low:app-B:eq:b-6}
\]
Thus the blue-minus-red gap is \(5/12+o(1)\), proving locally that red is
the active constraint for sufficiently large fixed \(C\).

\subsection{Order of quantifiers and proof status}

The absolute constant \(K\) in (2.10) is fixed before selecting the
large-\(C\) threshold, but its numerical value is not provided by the pinned
source.  The theorem therefore has the order
\[
 \exists C_0\ \forall C\ge C_0\ \exists\ell_0(C)\
 \forall\ell\ge\ell_0(C),
\tag{L-B.B.7}\label{low:app-B:eq:b-7}
\]
followed by \(\ell\to\infty\).  Only after that fixed-\(C\) theorem is
established is \(C\to\infty\) used to obtain the coefficient \(1/64\).
No explicit value of \(C_0\) or \(\ell_0(C)\) is asserted.

The resulting theorem is source-relative: its local residual algebra,
three-factor bridge, reverse-induction charge, extraction sum, first moments
and red--blue crossing are proved here.  Items (S1), (S2), (S4), and (S5)
are version-pinned source hypotheses; (S3) is the separately stated weighted
extension rule and is not attributed verbatim to either source.  The
arithmetic checker corroborates the finite formulas but does not replace
these probabilistic interfaces.
